%% kapa: exam=1 round=35 year=2024 period=1 form=A subject=civil_law subject_ko=민법 seq=26 number=26
%% source: Q-Net 감정평가사 기출문제 (cst00309) artlSeq=5213229 file=raw/1008/5213229/1교시 문제지 원본.pdf
%% conversion: kordoc 4.24.51 -> tools/gen1.py (auto); review=proofread-sonnet; flags=none
\begin{question}{26}
상린관계에 관한 설명으로 옳지 않은 것은?
\begin{choices}
  \item 경계에 설치된 담이 공유인 경우, 공유자는 그 분할을 청구할 수 있다.
  \item 인접하여 토지를 소유한 자는 다른 관습이 없으면 공동비용으로 통상의 경계표나 담을 설치할 수 있다.
  \item 경계표 설치를 위한 측량비용은 다른 관습이 없으면 토지의 면적에 비례하여 부담한다.
  \item 인접지의 수목뿌리가 경계를 넘은 경우, 토지소유자는 임의로 그 뿌리를 제거할 수 있다.
  \item 건물을 축조함에는 특별한 관습 또는 약정이 없으면 경계로부터 반미터 이상의 거리를 두어야 한다.
\end{choices}
\end{question}
\answer{1}
